\documentclass[12pt,a4paper]{article}

% ---------------------------------------------------------
% Codificação e idioma
% ---------------------------------------------------------
\usepackage[utf8]{inputenc}
\usepackage[T1]{fontenc}
\usepackage[brazil]{babel}

% ---------------------------------------------------------
% Layout e formatação
% ---------------------------------------------------------
\usepackage{geometry}
\usepackage{enumitem}
\usepackage{array}
\usepackage{amsmath}

\geometry{
    top=2.5cm,
    bottom=2cm,
    left=2cm,
    right=2cm,
    footskip=0.8cm
}

% ---------------------------------------------------------
% Imagens
% ---------------------------------------------------------
\usepackage{graphicx}

% ---------------------------------------------------------
% Código-fonte
% ---------------------------------------------------------
\usepackage{listings}
\usepackage{xcolor}

\lstset{
    language=C,
    basicstyle=\ttfamily\small,
    keywordstyle=\bfseries,
    commentstyle=\itshape,
    numbers=left,
    numberstyle=\tiny,
    numbersep=8pt,
    frame=single,
    breaklines=true,
    showstringspaces=false,
    tabsize=4
}

% ---------------------------------------------------------
% Cabeçalho e rodapé
% ---------------------------------------------------------
\usepackage{fancyhdr}
\usepackage{lastpage}

\pagestyle{fancy}
\fancyhf{}

% Cabeçalho esquerdo
\lhead{
    \begin{tabular}{l}
        \textbf{Estruturas de Dados - Avaliação}
    \end{tabular}
}

% Logo no canto superior direito
\rhead{
    \includegraphics[height=1.3cm]{logo.png}
}

% Rodapé
\cfoot{Página \thepage\ de \pageref{LastPage}}

% Altura necessária para acomodar o logo
\setlength{\headheight}{40pt}

% ---------------------------------------------------------
% Formatação dos parágrafos
% ---------------------------------------------------------
\setlength{\parindent}{0pt}
\setlength{\parskip}{6pt}

% ---------------------------------------------------------
% Documento
% ---------------------------------------------------------

\begin{document}

% =========================================================
% Identificação
% =========================================================

\begin{center}
    {\Large \textbf{Avaliação -- Estruturas de Dados}}\\[0.2cm]
    {\large Tecnologia em Análise e Desenvolvimento de Sistemas}
\end{center}

\vspace{0.5cm}

\noindent
\textbf{Nome:} \rule{7cm}{0.4pt}
\hfill
\textbf{RA:} \rule{3cm}{0.4pt}
\hfill
\textbf{Data:} 28/09/2026

\vspace{0.5cm}

\textbf{Orientações:}

\begin{itemize}
    \item Leia atentamente todas as questões antes de respondê-las.
    \item Considere a linguagem de programação C em todas as questões.
    \item Nas questões de programação, serão considerados o raciocínio,
          a estrutura da solução e o uso correto dos conceitos estudados.
    \item Respostas dissertativas devem ser justificadas de forma clara.
    \item Questões deixadas a lápis não poderão ser revistas.
    \item Questões com má caligrafia, serão desconsideradas.
\end{itemize}
\rule{16.5cm}{0.4pt}
\hfill
\subsection*{Questão 01}

Considere o código a seguir:

\begin{lstlisting}
int valores[5] = {10, 20, 30, 40, 50};
int *p = valores;

p = p + 2;

printf("%d", *p);
\end{lstlisting}

Assinale a alternativa que apresenta corretamente o valor exibido pelo
programa.

\begin{enumerate}[label=\Alph*)]
    \item 10
    \item 20
    \item 30
    \item 40
    \item 50
\end{enumerate}

\vspace{0.5cm}

% ---------------------------------------------------------

\subsection*{Questão 02}

Um programa precisa abrir o arquivo \texttt{alunos.txt} para
adicionar novos registros ao final do arquivo, preservando todo
o conteúdo que já existe.

Qual das alternativas realiza corretamente essa operação?

\begin{enumerate}[label=\Alph*)]
    \item \texttt{fopen("alunos.txt", "r");}
    \item \texttt{fopen("alunos.txt", "w");}
    \item \texttt{fopen("alunos.txt", "a");}
    \item \texttt{fopen("alunos.txt", "rw");}
    \item \texttt{fopen("alunos.txt", "n");}
\end{enumerate}

\subsection*{Questão 03}

Analise cuidadosamente o programa a seguir:

\begin{lstlisting}
#include <stdio.h>

void alterar(int *a, int b) {
    *a = *a + b;
    b = b * 2;
}

int main() {
    int x = 5;
    int y = 3;

    alterar(&x, y);

    printf("x = %d\n", x);
    printf("y = %d\n", y);

    return 0;
}
\end{lstlisting}

Responda:

\begin{enumerate}[label=\alph*)]
    \item Qual será o valor apresentado para \texttt{x}?   
    \item Qual será o valor apresentado para \texttt{y}? 
    \item Explique por que a função \texttt{alterar()} consegue modificar
    o valor de \texttt{x}, mas a alteração realizada em \texttt{b} não
    modifica \texttt{y}.
\end{enumerate}

% ---------------------------------------------------------

\subsection*{Questão 04}

Analise o programa:

\begin{lstlisting}
#include <stdio.h>

int main() {
    FILE *arquivo;
    char linha[100];

    arquivo = fopen("alunos.txt", "r");

    fgets(linha, sizeof(linha), arquivo);
    printf("%s", linha);

    fclose(arquivo);

    return 0;
}
\end{lstlisting}

O programador deseja ler todos os números existentes no arquivo
\texttt{numeros.txt}.

Entretanto, o código possui dois problemas importantes

\begin{enumerate}[label=\alph*)]
    \item Identifique os dois problemas.
    \item Explique quais consequências eles podem causar durante a execução.
    \item Apresente uma forma correta de realizar essa leitura.
\end{enumerate}

\subsection*{Questão 05 -- Ponteiros}

Explique, com suas próprias palavras, o que é um ponteiro na
linguagem C.

Em sua resposta, diferencie:

\begin{itemize}
    \item uma variável comum;
    \item o endereço de uma variável;
    \item uma variável ponteiro;
    \item os operadores \texttt{\&} e \texttt{*}.
\end{itemize}

Utilize um pequeno exemplo para complementar sua explicação.

% ---------------------------------------------------------

\subsection*{Questão 06 -- Alocação Dinâmica}

Durante a execução de um programa, nem sempre é possível conhecer
antecipadamente a quantidade de dados que será necessário armazenar.

Explique o que é alocação dinâmica de memória e por que ela
pode ser útil nessa situação.

Em sua resposta, explique também a finalidade das funções:

\begin{itemize}
    \item \texttt{malloc()};
    \item \texttt{calloc()};
    \item \texttt{realloc()};
    \item \texttt{free()}.
\end{itemize}
% ---------------------------------------------------------

\subsection*{Questão 07 -- Structs}

Considere um sistema que precisa armazenar as seguintes informações
sobre um aluno:

\begin{itemize}
    \item RA;
    \item nome;
    \item nota;
    \item situação acadêmica.
\end{itemize}

Explique por que uma \texttt{struct} é adequada para representar um
aluno nesse sistema.

Em seguida, explique a diferença entre utilizar uma \texttt{struct}
para representar um aluno e utilizar um vetor de
structs para armazenar vários alunos.

Não é necessário implementar o programa.

% ---------------------------------------------------------

\subsection*{Questão 08 -- Arquivos}

Explique por que um programa pode utilizar arquivos para
armazenar informações em vez de manter todos os dados exclusivamente
em variáveis na memória.

Em sua resposta, explique:

\begin{itemize}
    \item o que acontece com as variáveis quando o programa é encerrado;
    \item o significado de persistência de dados;
    \item para que serve o tipo \texttt{FILE *};
    \item a função de \texttt{fopen()} e \texttt{fclose()};
    \item a diferença entre os modos de abertura \texttt{"r"},
          \texttt{"w"} e \texttt{"a"}.
\end{itemize}

Utilize como exemplo um sistema que precisa manter informações
cadastradas mesmo depois que o programa for encerrado.

\subsection*{Questão 09}

Uma escola deseja armazenar temporariamente as notas de uma turma.

Crie um programa em C que:

\begin{enumerate}
    \item solicite ao usuário a quantidade de alunos;
    \item crie dinamicamente um vetor de \texttt{float} com espaço
          suficiente para armazenar as notas;
    \item leia todas as notas;
    \item calcule e apresente a média da turma;
    \item libere corretamente a memória utilizada antes do encerramento
          do programa.
\end{enumerate}

% ---------------------------------------------------------

\subsection*{Questão 10}

Considere a seguinte estrutura:

\begin{lstlisting}
typedef struct _produto{
    int codigo;
    char nome[50];
    float preco;
} Produto;
\end{lstlisting}

Escreva um programa ou trecho de código em C que:

\begin{enumerate}
    \item declare uma variável do tipo \texttt{Produto};
    \item leia do usuário o código, o nome e o preço do produto;
    \item abra o arquivo de texto \texttt{produtos.txt};
    \item acrescente ao final do arquivo os dados do produto;
    \item feche corretamente o arquivo.
\end{enumerate}

Os dados devem ser armazenados em uma única linha utilizando o seguinte
formato:

\begin{verbatim}
codigo;nome;preco
\end{verbatim}

Exemplo de conteúdo esperado no arquivo:

\begin{verbatim}
15;Teclado USB;89.90
\end{verbatim}

% =========================================================

\end{document}